\section{Notwendigkeit}

Änderungen im Verfassungsgefüge sind notwendig:
\begin{enumerate}[label=\alph*)]
    \item zur Härtung der Demokratie
    \begin{itemize}
        \item gegen Populismus
        \item gegen Kriminalität und Korruption in gewählten Ämtern
        \item gegen demokratisch gewählte Demokratiefeinde
        \item gegen Unfähigkeit der Spitzenpolitiker vor allem des Bundeskanzlers
        \item gegen folgenschwere falsche Entscheidungen von Exekutive und Legislative
    \end{itemize}
    \item zur Anpassung an das vorgeschlagene neue Wahlsystem
\end{enumerate}

\section{Härtung der Demokratie gegen Populismus}
\subsection{Motiviation}

Die Auseinandersetzung mit dem Populismus muss auf mehreren Ebenen erfolgen.
Wirklich besiegen kann man ihn nur, wenn man ihm die inhaltlichen Grundlagen entzieht.

Offenbar ist dies aber nicht so einfach, denn auch in der Demokratie gibt es Elitenbildung und viel schlimmer noch: Abkopplung der Eliten von der Lebensrealität der Nicht-Eliten.
Die Existenz von Eliten an sich ist eher etwas Gutes: Wir alle wollen nicht von Menschen geführt werden, die dümmer sind als wir.
Eigentlich wollen wir auch nicht von Menschen geführt werden, die genauso klug sind wie wir.
Stattdessen wollen wir von den Besten unter uns geführt werden.

Problematisch wird es, wenn immer größere Teile der Bevölkerung den Eindruck gewinnen, dass
\begin{enumerate}[label=\alph*)]
    \item die Eliten in ihrer eigenen Welt leben, die mit der Realität nichts mehr zu tun hat, oder
    \item nur ihre eigenen Interessen verfolgen, die den Interessen des \enquote{Normalbürgers} entgegenlaufen, oder
    \item die Zeiten nicht richtig einschätzen können, oder
    \item einfach auch keine Lösung haben, oder
    \item die Eliten gar nicht wirklich klüger sind als sie, sondern nur privilegiert.
\end{enumerate}

Dem kann man eigentlich nur entgegenwirken, indem man neu Menschen von unten nach oben bringt.
Die Demokratie lässt dies im Gegensatz zur Monarchie ja auch ohne Revolution zu.

Allerdings macht sie es auch nicht unbedingt einfach, von unten nach oben aufzusteigen.
Mit wohlhabenden Eltern in der richtigen Partei hat man schon deutlich bessere Aufstiegschancen als ein Kind aus bildungs- und politikfernen Schichten.

Trotzdem ist es möglich, dass auch Menschen, die das Leben kennen, bis an die Spitze kommen.
Es mag aber dauern, bis es wieder jemand schafft, in eine Führungsposition zu kommen, der wirklich eine gute Führungspersönlichkeit ist, und nicht nur jemand, der gut reden kann oder die notwendige Portion Vitamin B hat.



\section[Kontrolle von Parlament und Regierung]{Kontrolle von Parlament und Regierung durch nicht-gewählte Verfassungsorgane}

Es liegt in der Natur der Sache, dass vom Bürger gewählte Verfassungsorgane besonders anfällig gegen Populismus sind.
Auch nicht vom Bürger gewählte Verfassungsorgane können mit der Zeit unterwandert werden, indem sie nach und nach mit Demokratiefeinden besetzt werden.
Solche Prozesse brauchen aber Zeit.
Deswegen können nicht-gewählte Verfassungsorgane eine Zeit lang dem Populismus Grenzen setzen.

Der Gedanke ist, den Bürger die gestaltenden Verfassungsorgane (Exekutive und Legislative) selbst wählen zu lassen.
Deren Arbeit aber von nicht-gewählten Verfassungsorganen kontrollieren zu lassen.
Die kontrollierenden Verfassungsorgane dürfen selbst nicht gestalten, wie das Land sich entwickelt.
Sie sind ja auch nicht gewählt worden.
Sie sollen aber kontrollieren, dass die gewählten Gestalter ihre Arbeit richtig machen.

Ein kontrollierendes Verfassungsorgan haben wir schon.
Das Bundesverfassungsgericht führt Normenkontrollen durch.
Ihm fehlen aber Durchsetzungsmittel, und es fällt keine staatspolitischen Entscheidungen, sondern stellt nur fest, was recht ist.

Ich schlage vor, das Bundespräsidialamt personell und inhaltlich zu einem Bundespräsidialrat (oder \enquote{Staatsrat}?) zu erweitern.
Der Bundespräsident als Staatsoberhaupt bekommt sechs \enquote{Staatsräte} oder \enquote{Bundespräsidialräte} zur Seite, mit denen er gleichberechtigt Entscheidungen trifft.
Der Bundespräsident leitet die Sitzungen und bestimmt, was auf die Tagesordnung kommt.

Der Bundespräsident behält seine repräsentierende Funktion.
Die Staatsoberhauptsfunktion soll aber auf sieben Menschen aufgeteilt werden, um der erhöhten staatspolitischen Verantwortung gerecht zu werden, die dieses Amt übernehmen soll.


\section{Ablöseregelungen}\label{sec:abloseregelungen}

Aus vielen Gründen kann es vorkommen, dass die ganze Regierung (Legislative und Executive) oder auch nur Teile davon
vorzeitig abgelöst werden müssen.
Damit das Land zu keinem Zeitpunkt ohne Regierung bleibt, muss es für alle denkbaren Fälle Ablöseregelungen geben.
Die aktuell gültigen Regeln können nicht unmodifiziert übernommen werden.
Dazu sind die vorgeschlagenen Verfahrensänderungen zu grundlegend.
Zum Beispiel soll der Bundeskanzler jetzt ja indirekt vom Volk gewählt werden und nicht mehr vom Bundestag.
Daher kann der Bundestag nicht einfach so einen neuen Bundeskanzler wählen.
Das käme nur als befristete Übergangslösung in Betracht.

Im Folgenden werden daher die denkbaren Fälle einzeln durchgegangen:

\subsection{Ablösung des Bundeskanzlers ohne Legislativänderung}\label{subsec:ablosung-des-bundeskanzlers-ohne-legislativanderung}
- Kanzlerneuwahlen
\subsection{Zerbruch einer Koalition mit anschließender Bildung einer neuen Koalition}
Neuwahl des Bundeskanzlers durch den Bundestag nach Methode
\begin{enumerate}[label=\alph*)]
  \item Derjenige Kandidat der stärksten Koalitionspartei, der die meisten Stimmen erhalten hat, wird Bundeskanzler.
  \item Derjenige Kandidat von allen Koalitionsparteien, der die meisten Stimmen erhalten hat, wird Bundeskanzler.
  \item Der alte Bundeskanzler bleibt
\end{enumerate}
Wenn keine Wahl erfolgreich: Neuwahlen

\subsection{Zerbruch einer Koalition ohne Bildung einer neuen Koalition}
- Minderheitsregierung
- Neuwahlen

\subsection{Vorzeitige Ablösung einer Einparteienregierung}
- Neue Koalition
- Neuwahlen

\subsection{Ablösung einzelner Minister}
- Bundeskanzler entlässt und ernennt neu

\subsection{Ablösung einzelner Parlamentarier}
- Listenkandidaten rücken nach
- Direktwahlkandidaten: Personennachwahl im zugewiesenen Wahlkreis


\section{Demokratisches Gegengewicht} \label{sec:demokratisches-gegengewicht}
Durch das vorgeschlagene Wahlverfahren bekommt die Regierung mehr Macht.
Dafür muss an anderer Stelle ein demokratischer Ausgleich geschaffen werden, damit die Gefahr des Festsetzens an der Macht wieder verringert wird.

Die Regierung hat mehr Macht bekommen, damit sie ihren Auftrag, zu gestalten, auch in herausfordernden Zeiten, wo viel entschieden werden muss, erfüllen kann.
Umso wichtiger ist es, dass es eine "demokratische Notbremse" gibt, mit der man eine schlechte Regierung auch vorzeitig wieder loswerden kann, wenn diese ihre Macht missbraucht oder durch Inkompetenz
verschwendet.
Es muss Mechanismen geben, eine vorzeitige Ablösung von Einzelpersonen oder der ganzen Regierung einzuleiten, wenn dies geboten ist.

Die Herausforderung ist, das Machtgleichgewicht so zu gestalten, dass eine Regierung einerseits genug Macht hat, um auch unpopuläre Entscheidungen durchzusetzen, und andererseits noch gestoppt werden kann,
wenn sie zu viel Schaden anrichtet.

Im Folgenden sollen daher verschiedene Fälle diskutiert werden:

\subsection[]{Einzelne Parlamentarier verstoßen gegen das Gesetz}

Einzelne Parlamentarier können bereits jetzt nach Art 46 GG strafrechtlich verfolgt werden, wenn der Bundestag dem zustimmt.
Darüberhinaus sollte der Bundestag das Recht haben, Abgeordnete für die Dauer des Verfahrens zu suspendieren und bei
nachgewiesenen schweren Vergehen das Mandat sogar rückwirkend zu entziehen, so dass der Abgeordnete alle Privilegien
verliert, die mit einem Bundestagsmandat verbunden sind.

Schwere Vergehen in diesem Sinne sind Vergehen
\begin{itemize}
  \item die sich vorsätzlich gegen die Bundesrepublik Deutschland richten, wie z. B. Spionage für eine fremde Macht, oder
  \item Vergehen, die nach allgemeinem Rechtsempfinden ein Volksvertreteramt ausschließen, wie z. B. Gewaltverbrechen oder Sexualstraftaten.
 \end{itemize}

99\% der Bevölkerung würden sich wohl nicht von einem Mörder oder Kinderschänder vertreten lassen wollen, auch wenn er anderweitig noch so brilliant ist.
Bei Steuerhinterziehung und Trunkenheitsfahrten hingegen könnten die Meinungen vielleicht geteilt sein.

Der Bundestag hat das Recht bei Nachweis eines schweren Vergehens, das Mandat ab sofort zu entziehen oder sogar rückwirken, wenn der Strafbestand schon die ganze
Legislaturperiode bestanden hat.
Der betroffene Abgeordnete hat das Recht, bei dem Bundesverfassungsgericht gegen diese Maßnahme zu klagen.
Gibt es umgekehrt im Bundestag keine Mehrheit für die Suspendierung oder den Mandatsentzug, können einzelne Abgeordnete beim Bundesverfassungsgericht darauf klagen, und das
Bundesverfassungsgericht kann der Klage entsprechen oder nicht.


 {abloesung_einzelner_parlamentarier}
